% ── SECTION 2: THE CASE AGAINST SUBSTANCE ONTOLOGY (~2,000 words) ────────────

\section{What Quantum Mechanics Requires: The Case Against Substance Ontology}
\label{sec:substance}

\subsection{Substance Ontology and Its Requirements}

Western metaphysics from Aristotle~\citeyearpar{Aristotle_Metaphysics} through
Newton assumed that reality consists of persistent things --- substances ---
that have intrinsic properties and undergo change.
The substance is primary; its relations and processes are secondary.
Things exist first; then they relate.

For substance ontology to be compatible with physics, substances must have
\emph{intrinsic properties}: properties they possess independently of their
relations to other things.
A billiard ball has mass, charge, and position.
These are properties of the ball, not of its relations.

\subsection{Step 1: Bell's Theorem Eliminates Local Intrinsic Properties}

\citet{Bell1964} proved that no theory assigning definite values to
observables prior to measurement --- no matter how complex or hidden ---
can reproduce the statistical predictions of quantum mechanics without
violating locality.
\citet{CHSH1969} cast Bell's argument into the experimentally testable
inequality form (the CHSH inequality) that subsequent experiments have
used.
Experimental tests~\citep{Aspect1982,Hensen2015,Giustina2015,Shalm2015}
have confirmed Bell inequality violations, culminating in fully loophole-free
precision with both locality and detector-efficiency loopholes simultaneously closed.
The experimental program was recognized by the 2022 Nobel Prize in Physics,
awarded jointly to Aspect, Clauser, and Zeilinger~\citep{Zeilinger2023Nobel}.
Quantum particles do not have determinate spin, position, or momentum
before measurement --- not because we fail to measure them, but because
those properties do not, in any local theory, exist prior to the relational event of measurement.

The force of Bell's result deserves emphasis.
It is not a claim about our ignorance.
It is not a claim that there are hidden \emph{local} properties we happen
not to know.
Bell's theorem, in conjunction with the experimental record, establishes
that there are no such local properties to be ignorant of.
The \citet{Einstein1935EPR} program --- the hypothesis that quantum mechanics
is incomplete and that definite \emph{local} values exist prior to measurement,
awaiting only a more complete theory to reveal them --- is formally closed.
A nonlocal theory in which definite values are dynamically determined by a
global wave function (the de Broglie--Bohm interpretation) remains formally
available; I address it directly in \S\ref{sec:objections} below and argue
that it does not restore substance ontology so much as formalize one version
of the relational reading.
Quantum particles do not have intrinsic properties in the classical,
locally-instantiated sense.
What they have is something different: a relational potential, which is
actualized only through interaction.

\subsection{Step 2: Formal Incompatibility with Substance Ontology}

The consequence for substance ontology follows directly.
A substance, in the classical philosophical sense, is an entity that
possesses intrinsic properties --- properties it holds independently of
its relations to other things.
Bell's theorem and its experimental confirmation have proven that no
\emph{local} theory can attribute such properties to quantum particles.
A quantum particle, considered in isolation from relational context, has
no definite position, spin, or momentum that is locally specifiable.
The only formally available escape --- the nonlocal pilot-wave theory of
de Broglie and Bohm --- preserves particle positions only by determining
them dynamically from a global, configuration-space wave function (see
\S\ref{sec:objections}); on either reading, the Aristotelian picture of
locally self-contained substances has no instantiation in quantum reality.

This is not a matter of quantum mechanics being strange or counterintuitive.
It is a formal incompatibility with \emph{local} substance ontology.
Substance ontology, in its classical local form, requires intrinsic
properties.
Quantum mechanics, as confirmed by six decades of experimental tests of
Bell's inequalities, prohibits them at the most fundamental level of
physical description we possess.
The conclusion is that local substance ontology is not an alternative
interpretation of quantum reality but an interpretation that quantum
reality has formally refuted; the surviving nonlocal substance-style
reading (Bohmian mechanics) preserves the label of ``substance'' only
by importing the relational structure substance ontology was originally
meant to avoid.

The point is worth stating with precision.
The argument is not that substance ontology is philosophically uninteresting
or that process ontology is aesthetically preferable.
The argument is logical: \emph{substance without intrinsic properties
is not substance in any recognizable sense}.
It is an empty placeholder --- a grammatical subject without predicates.
If the only entities at the bottom of physical reality are ones that possess
no intrinsic properties, then the metaphysical category of substance,
understood classically, has no instantiation in the physical world.

\paragraph{The conceptual upshot.}
For the non-physicist, the formal result translates as follows.
The classical picture of reality as composed of self-contained objects
with their own intrinsic properties --- objects that could in principle
be described in isolation from any relation they might enter into ---
has not survived the experimental closure of local realism.
There is no ``billiard ball'' at the bottom of the world: only a network
of relational events whose properties become definite through interaction.
The death of the independent object is the conceptual hook that connects
the physics to Whitehead.
Process metaphysics is, at root, the philosophical articulation of a
world in which no entity is independent of its relations --- and that
world is the world the experimental record now requires us to inhabit.

\subsection{Step 3: The Wave Function as Description of Potential Events}

If quantum particles are not substances, what are they?
The formalism itself suggests an answer.

The wave function $\psi$ is not a description of what a particle \emph{is}.
It is a description of what relational events are \emph{possible} for it ---
a complete specification of the probabilities of various outcomes, conditional
on various measurement interactions.
The wave function does not say: the particle is here, spinning this way.
It says: if this system interacts with that system in that manner, the
following outcomes are possible with the following probabilities.

This is not a deficiency of the formalism.
It is a feature.
The wave function is precisely calibrated to describe relational potential:
the set of possible events that can actualize when a quantum system enters
a measurement relationship.
Heisenberg's uncertainty principle~\citeyearpar{Heisenberg1927} is not,
on this reading, a limitation on measurement precision.
It is a statement about the ontological structure of quantum systems:
the properties that can be jointly definite are constrained, because
definiteness is not intrinsic but relational, and different measurement
relationships actualize different aspects of the system's potential.

In the vocabulary of process philosophy, the wave function describes an
actual occasion in its \emph{potential phase} --- a concrete quantum system
in the process of becoming, prior to the relational event that will actualize
a definite outcome.
This identification is not an analogy but a substantive interpretive
proposal, developed in detail by \citet{Epperson2004}: the wave function
is read as the formal description of what Whitehead's process ontology
would lead us to expect of a pre-actualization quantum state --- a
distribution of relational possibilities, not a catalogue of intrinsic
properties.
The mapping is not a proof that process ontology is the correct
interpretation of the formalism, but it does establish that the formalism
admits a process-ontological reading without distortion --- which is more
than can be said for any reading that requires intrinsic, locally
specifiable properties.

\subsection{Step 4: Relational Quantum Mechanics}

\citet{Rovelli1996} proposes that quantum states are always relative to an
observer or system --- never absolute.
On my reading, the measurement event in Rovelli's framework is the moment
when superposition resolves into a definite outcome through relational
interaction --- the structure of what process ontology calls an actual
occasion: a concrete event of relational becoming in which potential is
actualized through interaction.

There is no view from nowhere in Rovelli's relational quantum mechanics.
A particle does not have spin-up; it has spin-up \emph{relative to this
detector, in this measurement event}.
This is not subjectivism --- the outcome is real, and it is determinate.
But it is irreducibly relational: the definiteness of the outcome exists
only within the context of a specific relational event.
Two observers who have not compared notes may assign different states to
the same system --- not because one of them is wrong, but because quantum
states are facts about relations, and different observers stand in
different relational contexts~\citep{Rovelli2021}.

This is, on my reading, precisely the structure that process ontology predicts.
If reality is constituted by relational events rather than intrinsic
properties, then the quantum mechanical description --- in which states
are observer-relative and properties are actualized through interaction ---
is not a puzzle to be explained away.
It is what physics looks like when the underlying metaphysics is correct.

\subsection{Surviving Interpretations and Replies}
\label{sec:objections}

The claim that quantum mechanics requires process ontology will be resisted.
Four families of interpretation present apparent counterexamples and
deserve direct response.

\paragraph{Bohmian (pilot-wave) mechanics.}
The de Broglie--Bohm interpretation~\citep{Bohm1980} is the most direct
candidate for preserving substance ontology under the experimental record:
particles have definite trajectories at all times, guided by a pilot wave
that propagates nonlocally on the system's configuration space.
Of the surviving interpretations, Bohmian mechanics is the only one that
appears to restore the classical picture of particles with intrinsic
positions.

The reply is that Bohm restores positions only by importing the very
relational structure substance ontology was meant to avoid.
The pilot wave is a global object defined on configuration space, and the
trajectory of any single particle depends instantaneously on the positions
of every other particle in the system.
The ``intrinsic'' position of a Bohmian particle is dynamically determined,
at every instant, by a nonlocal field that encodes the relations among all
particles in the system.
This is substance ontology in name only: the supposedly intrinsic property
is in fact the output of a holistic, nonlocally-coupled guiding field.
Bohm does not, on my reading, refute the relational reading of quantum
mechanics; he formalizes, on my reading, one version of it, with the
additional commitment to point particles whose positions do no explanatory
work the wave function does not already do.
A reader who accepts Bohmian mechanics has, in effect, accepted that the
fundamental physical object is a nonlocal relational field --- which is
the substantive content of the process-ontological claim, dressed in
particle vocabulary.

\paragraph{Many-worlds interpretation.}
The Everett many-worlds interpretation~\citep{Everett1957,Wallace2012} avoids the
measurement problem by denying wave function collapse: the wave function
evolves unitarily, and what we experience as a measurement outcome is
our branch of a universal branching structure.
The objection is that many-worlds preserves a form of substance ontology:
the universal wave function is an absolute, observer-independent object.

The reply has two parts.
First, many-worlds resolves the measurement problem by proliferating
ontology on an extraordinary scale: every quantum interaction ramifies
into a new branch of the universal wave function, yielding an uncountably
infinite set of simultaneously real worlds.
This does not so much solve the problem of intrinsic properties as
relocate it: intrinsic properties are now assigned not to particles
but to branches of the multiverse, and the question of what selects
our branch remains unresolved.
Second, many-worlds does not recover what Bell's theorem eliminated.
Even in a many-worlds framework, no single particle in any single
branch has definite intrinsic properties prior to branching.
The relational structure of quantum measurement is preserved in
each branch.
Many-worlds is a response to the \emph{measurement problem}, not
a restoration of \emph{intrinsic properties}.
Bell's argument stands.

\paragraph{QBism.}
Quantum Bayesianism (QBism)~\citep{FuchsMerminSchack2014} treats quantum states as
agents' beliefs about future experiences, not descriptions of objective
reality.
The objection is that if quantum states are just beliefs, the ontological
question does not arise: there is no fact of the matter about what the
particle really is.

The reply is that QBism sidesteps rather than answers the ontological question.
Proponents (notably \citealt{FuchsMerminSchack2014}) deny that an external
ontology is required, arguing that the formalism's structural constraints ---
including reformulations of the Born rule via symmetric informationally
complete measurements --- are themselves the substantive content of the
theory.
Granting this for argument's sake, the question still arises in altered form:
what is the structure of the world such that these constraints govern
rational credence in a quantum context?
QBism answers this only by appeal to a background ontology, whether explicit
or implicit in the structural constraints themselves --- and that ontology,
if it is to be consistent with quantum mechanics, will be relational rather
than substantival.
QBism defers the ontological question without eliminating it.

\paragraph{Epistemic interpretations.}
A family of epistemic interpretations holds that quantum states describe
not the state of a system but our knowledge of it.
The wave function is a probability distribution over some underlying
classical reality that we are ignorant of.

This is the position that Bell's theorem and its successors directly refute.
The \citet{Einstein1935EPR} argument assumed that there is an underlying
classical reality, with definite properties, that quantum mechanics
incompletely describes.
Bell's theorem proved that no such reality, if it is local, can reproduce
the statistical predictions of quantum mechanics.
The Pusey--Barrett--Rudolph theorem~\citep{PBR2012} sharpens the constraint
further: under modest independence assumptions, the quantum state cannot be
a state of knowledge about an underlying ontic state with definite properties
shared by distinct quantum states; the wave function must itself be ontic.
Epistemic interpretations of the standard $\psi$-epistemic kind are not
a live option in the presence of experimental Bell inequality violations
and the PBR result.
The correlations that quantum mechanics predicts --- and that experiment
has confirmed --- cannot be explained by ignorance of pre-existing classical
properties.
They require that the properties themselves come into being through
relational interaction.

The conclusion of all four steps stands: every interpretation of quantum
mechanics that survives the experimental record is, on my reading, either
a process ontology (relational quantum mechanics, Whiteheadian readings)
or a relabeled version of one (Bohmian mechanics as a nonlocal relational
field; many-worlds as branching relational structure within each branch;
QBism as a relational agent-world coupling).
The question is no longer whether to prefer process ontology as an
interpretation, but what process ontology actually says --- and whether
any tradition has previously articulated its structure.
